%%%%%%%%%%%%%%%%%%%%%%%%%%%%%%%%%%%%
% template.tex
%
% Template for the KTHEEposter class.
%
% Original version: Mats Bengtsson, 28/5 2002
% Updated:
%   Mats Bengtsson April 2011 - July 2012
%   Manuel Olguín October 2018
%%%%%%%%%%%%%%%%%%%%%%%%%%%%%%%%%%%%

<<<<<<< HEAD
\documentclass[portrait, a0]{KTHEEposter}
=======
\documentclass[portrait, a1]{KTHEEposter}
>>>>>>> init writing
\usepackage{lipsum}
\usepackage{listings}
\usepackage{enumitem}
\usepackage{caption}
<<<<<<< HEAD
\lstset{basicstyle=\ttfamily, frame=single}

\kthlogo{kth_eng_cmyk}
%\extralogo{kth_eng_cmyk}
=======
\usepackage{indentfirst}
\lstset{basicstyle=\ttfamily, frame=single}

\kthlogo{kth_eng_cmyk}
\extralogo{img/cmu_csd_logo}
>>>>>>> init writing
\kthcolor{KTHblue}

\begin{document}
    
<<<<<<< HEAD
    \title{KTH Royal Institute of Technology\\Poster Template}
    
    \author{Mats Bengtsson, Manuel Olguín}
    \maketitle

    % begin poster columns:
    \begin{pcolumns}[3] % optional argument is the number of columns, by default it's 2
        \begin{pcolumn}[4] % optional argument is the number of frames in the column, by default 2
            \begin{pframe}[1.15] % optional argument indicates the scaling of the frame size. Sum of this for all frames in a column should be equal to the number of frames!
                The KTHEEposter class provides helper functions to easily set up posters which follow the KTH visual style.
                In particular it provides the following environments:
                \begin{itemize}
                    \item To initialize a column environment with pretty frames:
                    \begin{lstlisting}
\begin{pcolumns}[<# of cols>]
... 
\end{pcolumns}
                    \end{lstlisting}
                    
                    \item To set up a column inside a column environment:
                    \begin{lstlisting}
\begin{pcolumn}[<# of frames>]
...
\end{pcolumn}
                    \end{lstlisting}
                
            
                    \item And then environments to set up frames\dots
                \end{itemize}
            \end{pframe}
            \begin{pframe}[.35] % frame with border and shadow
                \dots \textbf{with border and shadow:}
                \begin{lstlisting}
\begin{pframe}[<scale>]
...
\end{pframe}
                \end{lstlisting}
            \end{pframe}
            \begin{pframe*}[.35] % frame with border, no shadow
                \dots \textbf{with border but no shadow:}
                \begin{lstlisting}
\begin{pframe*}[<scale>]
...
\end{pframe*}
                \end{lstlisting}
            \end{pframe*}
            \begin{pframe**}[.35] % frame with no border, no shadow
                \dots \textbf{no border and no shadow:}
                \begin{lstlisting}
\begin{pframe**}[<scale>]
...
\end{pframe**}
                \end{lstlisting}
            \end{pframe**}
        \end{pcolumn}%
        \begin{pcolumn}[3]
            \begin{pframe}[.7]
                The class also provides easy options for changing the colors of the frames.
                Simply invoke the \texttt{\textbackslash kthcolor\{\}} command in your preamble with one of two color options:
                \begin{itemize}
                    \item \texttt{KTHblue} for the standard KTH blue color on the frames.
                    \item \texttt{KTHred} for the alternative KTH red color.
                \end{itemize}
                \medskip
                If you want to insert an alternative logo instead of the default KTH logo on the right side of the title, use the \texttt{\textbackslash extralogo\{\}} command, passing the desired image as a parameter.
            \end{pframe}
            \begin{pframe}[.2]
                See the source code for this template (\texttt{template.tex}) and for the class itself (\texttt{KTHEEposter.cls}) for more details.
            \end{pframe}
            \begin{pframe}[2.1]
                Of course, you can add arbitrary \LaTeX~code inside the frames, with the exception of floats (figures, etc).
                One solution is to use the \texttt{caption} package to add correct captions inside a minipage (see source code):
                
                \vskip 20mm
                
                \begin{minipage}{\linewidth}
                    \centering
                    \includegraphics[width=.7\linewidth]{Logo/kth_eng_cmyk}
                    \captionof{figure}{KTH Logo}
                \end{minipage}
            \end{pframe}
        \end{pcolumn}%
        \begin{pcolumn}[3]
            \begin{pframe**}[1]
            \end{pframe**}
            \begin{pframe}[.5]
                The KTHEEposter document class was originally created by Mats Bengtsson in 2002, who also updated it in 2014.
                This new version by Manuel Olguín, created in 2018, builds upon Mats' implementation and modernizes it a bit.
                Some functionality of the old class has not yet been implemented though, in particular the arrows between frames.
                This will be re-added in the near future!
            \end{pframe}
            \begin{pframe}[1.5]
                \section{Changelog}
                \begin{description}
                    \item[8 October 2018] Reimplement using tcolorbox. 
                    Add options to change colors. 
                    General reformatting of the class. M. Olguín.
                    \item[2 April 2014] Add latest version of the KTH logo. M. Bengtsson.
                    \item[July 2012] Renewed layout, reimplemented using tikz instead of pstricks, added possibility to add extra logotype. M. Bengtsson.
                    \item[7 August 2006] Renamed to KTHEEposter and modified to adhere somewhat more to the graphical profile. M. Bengtsson.
                    \item[28 May 2002] Original version. M. Bengtsson.
                \end{description}
            \end{pframe}
        \end{pcolumn}%
    \end{pcolumns}%
=======
    \title{\LARGE\bfseries Scaling on the Edge:\\A Benchmarking Suite For Human-in-the-Loop Applications}
    
    \author{\Large M. Olguín, J. Wang, M. Satyanarayanan, J. Gross}
    \maketitle
    
    \begin{pcolumns}[3]
        \begin{pcolumn}[2]
            \begin{pframe}
                Benchmarking human-in-the-loop application is complex given their nature, which heavily depends on the actions taken by the \emph{human} user.
                This limits reproducibility as well as feasibility of performance evaluations.
                In this paper we propose a methodology and present a benchmarking suite we call EdgeDroid  that can address these challenges.
                Our core idea rests on recording traces of these applications which are played out in a controlled fashion based on an underlying model of human behavior.
                The traces are then exposed to the original backend compute process of the respective human-in-the-loop application, generating realistic feedback.
                This allows for an automated system which greatly simplifies benchmarking large scale scenarios.
                After presenting this methodology and implementation, we demonstrate the approach on an example cognitive assistance application.
                The results show the utility of EdgeDroid as a tool for system designers, application developers and researchers alike.
            \end{pframe}
        \end{pcolumn}
    \end{pcolumns}
>>>>>>> init writing
    
\end{document}